\documentclass[11pt]{article}
\usepackage{mathunal-base}
\mathunalfuente{serif}

\renewcommand{\examtitulo}{Ecuaciones Diferenciales (1000007) --- Segundo Examen Parcial}
\renewcommand{\examfecha}{Semestre 02-2015, 17 de Octubre de 2015, TEMA A}

\begin{document}

\examheader

\vspace{4pt}
\noindent
\begin{minipage}[t]{0.62\linewidth}
\textbf{Duración: 1 hora y 40 minutos}

\vspace{8pt}
Nombre y carnet: \dotfill

\vspace{6pt}
Profesor y grupo: \dotfill
\end{minipage}%
\hfill
\begin{minipage}[t]{0.32\linewidth}
\begin{tabular}{|l|c|}
\hline
\multicolumn{2}{|c|}{\textbf{Calificación}} \\ \hline
1. & \\ \hline
2. & \\ \hline
3. & \\ \hline
4. & \\ \hline
\textbf{Total} & \\ \hline
\end{tabular}
\end{minipage}

\vspace{6pt}
\noindent\textbf{Instrucciones.} No está permitido el uso de calculadora ni de cualquier otro tipo de aparato electrónico, incluyendo teléfonos móviles, los cuales deben permanecer apagados durante el tiempo de duración del examen. Tampoco está permitido hacer preguntas a las personas encargadas de la vigilancia.

\vspace{10pt}
\begin{enumerate}
\item{} (30\%) En los numerales i, ii, iii, iv y v complete según corresponda.
\begin{enumerate}
\item{} (6\%) Considere la ecuación diferencial
\[
2y'''+7y''+8y'+ky=0.
\]
\noindent El valor de $k$ para que $e^{-t}$ sea solución de la ecuación diferencial es \dotfill.

\item{} (6\%) Una forma adecuada de solución particular $y_p$ para la ecuación diferencial
\[
y^{(iv)}+2y'''+2y'' = 3e^x+2xe^{-x}+e^{-x}\operatorname{sen}(x)
\]
\noindent es \dotfill. (No evalúe constantes)

\item{} (6\%) Si una ecuación diferencial homogénea con coeficientes constantes tiene por solución general
\[
y(x) = e^{-2x}\left(c_1\cos(3x)+c_2\operatorname{sen}(3x)\right),
\]
\noindent entonces la ecuación diferencial es \dotfill.

\item{} (6\%) Los valores de la constante $a$ para los cuales las soluciones del problema de valor inicial
\[
y''-y=0, \qquad y(0)=-a,\quad y'(0)=2
\]
\noindent son siempre positivas son \dotfill.

\item{} (6\%) Una ecuación diferencial de Cauchy-Euler $ax^2y''+bxy'+cy=0$, donde $a,b,c$ son constantes reales, tiene a $r_1=-3$ y $r_2=1$ como raíces de la ecuación auxiliar. La ecuación diferencial es \dotfill.
\end{enumerate}

\item{} (20\%) Halle la solución general de ecuación diferencial
\[
(1-2x)^2y''(x)-(2x-1)y'(x)+2y(x)=0, \qquad 2x>1.
\]
\noindent (Sugerencia: Busque soluciones de la forma $y(x)=(2x-1)^m$ ó haga la sustitución $t=2x-1$)

\vspace{6cm}

\newpage
\item{} (30\%) Una masa que pesa 2 lb estira un resorte 12 pulgadas. A este resorte se le sujeta una masa que pesa 3.2 lb. El sistema masa-resorte se sumerge en un medio con constante de amortiguamiento 0.4 lb-seg/pie. La masa se pone en movimiento desde su posición de equilibrio con una velocidad de 3 pul/seg hacia abajo.
\begin{enumerate}
\item{} (15\%) Determine la posición $x(t)$ de la masa en cualquier instante $t$
\item{} (5\%) Bosqueje la gráfica de la posición $x(t)$.
\item{} (5\%) Halle el cuasi-período.
\item{} (5\%) Determine la segunda vez en que la masa pasa por la posición de equilibrio moviéndose hacia arriba.
\end{enumerate}
\noindent (Recuerde que la constante $g$ es 32 pies sobre segundo al cuadrado)

\vspace{6cm}

\item{} (20\%) Resuelva la ecuación diferencial
\[
y''+9y = \sec(3x), \qquad -\frac{\pi}{2}<x<\frac{\pi}{2}
\]

\vspace{6cm}
\end{enumerate}

\examfooter
\end{document}
